\documentclass[../../main.tex]{subfiles}
\begin{document}

% 年度の大きな表示
\begin{center}
  {\fontsize{48pt}{58pt}\selectfont \textbf{2002}\normalsize\textbf{年度}}
\end{center}

\vspace{10mm}

% 本文


\vspace{15mm}

% 出題テーマと難易度の表
\noindent\textbf{出題テーマと難易度}

\vspace{5mm}

\begin{center}
  \shadowbox{%
    \begin{tabular}{|c|p{8cm}|c|}
      \hline
      \textbf{問題番号} & \textbf{テーマ} & \textbf{難易度} \\
      \hline
      第1問 & 絶対値付き定積分の最小値（三角関数） & \\
      \hline
      第2問 & 楕円の直交する2接線と点の軌跡（準円） & \\
      \hline
      第3問 & 空間図形（正三角形の影の面積の最大化） & \\
      \hline
      第4問 & 調和数列の対数評価と方程式の解の極限 & \\
      \hline
    \end{tabular}%
  }
\end{center}

\vspace{15mm}

% 解答
\noindent\textbf{解答}

\vspace{5mm}

\begin{center}
  \shadowbox{%
    \renewcommand{\arraystretch}{1.6}
    \begin{tabular}{|c|c|p{8cm}|}
      \hline
      \textbf{問題番号} & \textbf{小問} & \textbf{答え} \\
      \hline
      第1問 & - & 最小値 $\dfrac{\sqrt{14}}{2}-\dfrac{\sqrt{2}}{2}-1$ \\
      \hline
      第2問 & - & 軌跡は円 $a^2+b^2=25$ \\
      \hline
      第3問 & - & 最大値 $\dfrac{(\sqrt{2}+1)^2}{2}$ \\
      \hline
      \multirow{3}{*}{第4問} & (1) & $1$ \\
      \cline{2-3}
                             & (2) & 証明問題 \\
      \cline{2-3}
                             & (3) & $1$ \\
      \hline
    \end{tabular}%
    \renewcommand{\arraystretch}{1.0}
  }
\end{center}

\vspace{15mm}

% 解答時間
\noindent\textbf{試験形態}

\vspace{5mm}

\begin{center}
  \shadowbox{%
    \begin{tabular}{p{8cm}}
      （要確認）
    \end{tabular}%
  }
\end{center}

\end{document}
